% Generated draft: requires mathematical and source review.
\documentclass[11pt]{article}
\usepackage[margin=25mm]{geometry}
\usepackage{fontspec}
\setmainfont{Latin Modern Roman}
\usepackage{amsmath,amssymb,mathtools,unicode-math}
\setmathfont{Latin Modern Math}
\usepackage{enumitem,needspace,xurl,hyperref}
\setlength{\parindent}{0pt}
\setlength{\parskip}{5pt}
\setlength{\emergencystretch}{3em}
\newcommand{\N}{\mathbb{N}}
\newcommand{\Z}{\mathbb{Z}}
\newcommand{\Q}{\mathbb{Q}}
\newcommand{\R}{\mathbb{R}}
\newcommand{\C}{\mathbb{C}}
\newcommand{\F}{\mathbb{F}}
\newcommand{\A}{\mathbb{A}}
\newcommand{\PP}{\mathbb{P}}
\newcommand{\E}{\mathbb{E}}
\newcommand{\eps}{\varepsilon}
\newcommand{\dd}{\,\mathrm{d}}
\begin{document}
\section*{UnsolvedMath ID: 701262}
\textbf{Source problem code:} catalogue entry
\subsection{Definitions and mathematical statement}
Let \(A\) be an ordered abelian group, with a translation-invariant total order. An \(A\)-tree is a nonempty \(A\)-valued metric space in which each pair of points is joined by a unique geodesic segment isometric to an interval in \(A\), and every geodesic triangle is a tripod. An action on such a tree is understood to be by isometries. We use the standard convention that a free action has trivial point stabilizers and has no inversions: no group element interchanges the endpoints of a nondegenerate segment that it preserves.

For \(n\geq 1\), equip \(\Z^n\) with lexicographic order: for distinct vectors \(u,v\), one has \(u<v\) if the first nonzero coordinate of \(v-u\) is positive. The supplied record does not explicitly specify the order on \(\Z^n\), or its convention concerning inversions; the formulation below uses these standard conventions and requires confirmation against the source.

\textbf{Question.} For every ordered abelian group \(A\) and every finitely generated group \(G\), does the following implication hold?
\[
\begin{gathered}
G\text{ admits a free isometric action on an }A\text{-tree}
\\[2pt]
\Longrightarrow
\\[2pt]
\exists n\in\Z_{\geq 1}\quad
G\text{ admits a free isometric action on a lexicographically ordered }\Z^n\text{-tree}.
\end{gathered}
\]
The integer \(n\) is not required to be uniform over the groups in the question.

\textbf{Supplied 2013 qualification.} The record reports that every finitely presented group acting freely on an \(A\)-tree acts freely on an \(\R^n\)-tree for some suitable \(n\), citing O. Kharlampovich, A. Myasnikov, and D. Serbin. Its accompanying order description literally refers to \(\R\) as having lexicographic order; the likely intended reference is lexicographic order on \(\R^n\), but that wording requires source confirmation. This reported result concerns finitely presented groups and an \(\R^n\)-tree target, whereas the question concerns finitely generated groups and a \(\Z^n\)-tree target. The report is supplied background, not an independently verified theorem or resolution of the question.
\subsection{Short English statement}
Must every finitely generated group that acts freely on a tree with distances in some ordered abelian group also act freely on a tree with distances in lexicographically ordered \(\Z^n\), for some positive integer \(n\)? The supplied record labels the problem open. Its 2013 comment reports a related result for finitely presented groups with an \(\R^n\)-tree target; that report and the order conventions require source review.
\subsection{Sources}
\begin{itemize}
\item Local UnsolvedMath registry: supplied complete source record, permanent ID \(701262\), exact source problem code \texttt{catalogue entry}, title \emph{Kourovka Notebook Problem 16.70}. The record supplies the status \texttt{open}; this status has not been independently verified. Its creation and update timestamps are both \texttt{2026-10-16T14:00:00Z}.
\item \emph{The Kourovka Notebook (Unsolved Problems in Group Theory)}, original problem 16.70, 16th issue (2006); the supplied record identifies the 21st issue edition (2026), \url{https://arxiv.org/abs/1401.0300v48}. Proposed by A. G. Myasnikov, V. N. Remeslennikov, and O. G. Kharlampovich. The collection is attributed to editors E. I. Khukhro and V. D. Mazurov. These bibliographic details and the link are supplied attribution, not independently checked source access.
\item Reference supplied in the comment of 2013: O. Kharlampovich, A. Myasnikov, D. Serbin, \emph{Int. J. Algebra Comput.} \textbf{23} (2013), 325--345. No article title or article-specific link is supplied. The associated mathematical claim is retained as a reported source qualification.
\end{itemize}
\end{document}
